\section{总结与体会}

本项目围绕量子线路在受限硬件拓扑上的路由编译展开，把量子信息基础中的线路模型、
双比特门、SWAP、硬件耦合图和线路深度落实到一个可运行系统中。报告先说明真实芯片
为什么需要路由，再复现 SABRE basic 作为传统启发式基线，随后给出 NPQR 的完整流程：
初始映射选择起点，前沿推进维护线路依赖，动作评分对候选 SWAP 排序，束搜索保留多条
路线，局部修复处理尾部停滞，复放验证保证输出线路满足拓扑约束。

量子信息基础大作业要求的理论、算法、实验和工程交付在项目中形成了闭环。理论部分覆盖量子态、
量子门、双比特门、SWAP、拓扑映射和路由代价；算法部分完成 SABRE 复现和 NPQR 设计；
实验部分给出 10/20 比特主实验、30/50 比特扩展实验、消融实验和阶段耗时分析；工程
部分实现网页演示、服务器接口、MCP 服务、命令行、本地测试、GitHub Pages 前端和
云端部署。AI 协作部分记录了资料整理、前端开发、服务接口、测试部署和报告写作中的
使用方式，算法目标、实验口径和结果解释由作者根据代码运行结果确定。

实验结果表明，NPQR 在 10/20 比特代表样例上完成 $10/10$，SWAP 数均低于 SABRE
basic；选定 30/50 比特网格样例完成 $4/4$，并保持 SWAP 优势。10 个主实验样例中，
NPQR 合计插入 218 个 SWAP，SABRE basic 合计插入 359 个，减少 141 个。阶段耗时
显示，不同线路的瓶颈并不相同：QAOA10 和 Brickwork20 主要受初始映射候选生成影响，
QFT10 的主搜索占总耗时 90.2\%。消融实验进一步说明，束搜索和结构化初始映射对完成率
和交换数都有直接影响。

项目的后续优化可以沿四个方向推进：改进大拓扑初始映射候选生成，压缩前沿相关候选
SWAP 集合，引入分层路由以降低搜索状态规模，并对模型推理做批处理。这些方向延续
当前设计原则：把学习模型放入显式、可复放的搜索框架中，用清楚的拓扑约束和实验指标
支撑量子线路路由编译的结果。
